
\documentclass{article}
\usepackage{colortbl}
\usepackage{makecell}
\usepackage{multirow}
\usepackage{supertabular}

\begin{document}

\newcounter{utterance}

\centering \large Interaction Transcript for game `bbh', experiment `logical\_deduction\_three\_objects', episode 68 with reflexion{-}llama{-}3.1{-}8b{-}instruct.
\vspace{24pt}

{ \footnotesize  \setcounter{utterance}{1}
\setlength{\tabcolsep}{0pt}
\begin{supertabular}{c@{$\;$}|p{.15\linewidth}@{}p{.15\linewidth}p{.15\linewidth}p{.15\linewidth}p{.15\linewidth}p{.15\linewidth}}
    \# & \multicolumn{2}{c}{Player} && \multicolumn{2}{c}{Game Master} \\
    \hline

    \theutterance \stepcounter{utterance}  
    & & & \multicolumn{4}{p{0.6\linewidth}}{
        \cellcolor[rgb]{0.9,0.9,0.9}{
            \makecell[{{p{\linewidth}}}]{
                \texttt{\tiny{[P1$\langle$GM]}}
                \texttt{A logical deduction task which requires deducing the order of a sequence of objects.} \\
\\ 
\texttt{Q: The following paragraphs each describe a set of three objects arranged in a fixed order. The statements are logically consistent within each paragraph. In a golf tournament, there were three golfers: Amy, Eli, and Eve. Eve finished above Amy. Eli finished below Amy.} \\
\texttt{Options:} \\
\texttt{(A) Amy finished last} \\
\texttt{(B) Eli finished last} \\
\texttt{(C) Eve finished last} \\
\texttt{A: (B)} \\
\\ 
\texttt{Q: The following paragraphs each describe a set of three objects arranged in a fixed order. The statements are logically consistent within each paragraph. On a shelf, there are three books: a white book, a green book, and an orange book. The green book is to the right of the white book. The orange book is the rightmost.} \\
\texttt{Options:} \\
\texttt{(A) The white book is the leftmost} \\
\texttt{(B) The green book is the leftmost} \\
\texttt{(C) The orange book is the leftmost} \\
\texttt{A: (A)} \\
\\ 
\texttt{Q: The following paragraphs each describe a set of three objects arranged in a fixed order. The statements are logically consistent within each paragraph. On a shelf, there are three books: a red book, a gray book, and a white book. The white book is to the left of the gray book. The red book is the second from the left.} \\
\texttt{Options:} \\
\texttt{(A) The red book is the leftmost} \\
\texttt{(B) The gray book is the leftmost} \\
\texttt{(C) The white book is the leftmost} \\
\texttt{A: (C)} \\
\\ 
\texttt{Q: The following paragraphs each describe a set of three objects arranged in a fixed order. The statements are logically consistent within each paragraph. On a branch, there are three birds: a cardinal, a hawk, and a hummingbird. The hummingbird is to the right of the hawk. The cardinal is to the left of the hawk.} \\
\texttt{Options:} \\
\texttt{(A) The cardinal is the rightmost} \\
\texttt{(B) The hawk is the rightmost} \\
\texttt{(C) The hummingbird is the rightmost} \\
\texttt{A:} \\
            }
        }
    }
     \\ \\

    \theutterance \stepcounter{utterance}  
    & \multicolumn{4}{p{0.6\linewidth}}{
        \cellcolor[rgb]{0.561,0.737,0.561}{
            \makecell[{{p{\linewidth}}}]{
                \texttt{\tiny{[P1$\rangle$GM]}}
                \texttt{To deduce the order of the birds, let's analyze the given information:} \\
\\ 
\texttt{1. The hummingbird is to the right of the hawk.} \\
\texttt{2. The cardinal is to the left of the hawk.} \\
\\ 
\texttt{From statement 2, we know the cardinal is to the left of the hawk. From statement 1, we know the hummingbird is to the right of the hawk. } \\
\\ 
\texttt{Combining these two statements, we can conclude that the order is: } \\
\texttt{cardinal {-}> hawk {-}> hummingbird} \\
\\ 
\texttt{So, the correct answer is:} \\
\texttt{(A) The cardinal is the rightmost} \\
            }
        }
    }
    & & \\ \\

    \theutterance \stepcounter{utterance}  
    & & & \multicolumn{2}{p{0.3\linewidth}}{
        \cellcolor[rgb]{0.9,0.9,0.9}{
            \makecell[{{p{\linewidth}}}]{
                \texttt{\tiny{[GM$|$GM]}}
                \texttt{Target: (C)} \\
            }
        }
    }
    & & \\ \\

    \theutterance \stepcounter{utterance}  
    & & & \multicolumn{2}{p{0.3\linewidth}}{
        \cellcolor[rgb]{0.9,0.9,0.9}{
            \makecell[{{p{\linewidth}}}]{
                \texttt{\tiny{[GM$|$GM]}}
                \texttt{game\_result = LOSE} \\
            }
        }
    }
    & & \\ \\

    \theutterance \stepcounter{utterance}  
    & & \multicolumn{4}{p{0.6\linewidth}}{
        \cellcolor[rgb]{0.95,0.95,0.95}{
            \makecell[{{p{\linewidth}}}]{
                \texttt{\tiny{[P1$\leftrightarrow$P]}}
                \texttt{Post{-}episode Reflexion (not sent to GM):} \\
\texttt{New plan: } \\
\\ 
\texttt{In my previous attempt, I correctly deduced the order of the birds as: cardinal {-}> hawk {-}> hummingbird. However, I incorrectly concluded that the cardinal is the rightmost bird. This mistake occurred because I only considered the two given statements and didn't think about the fact that the hummingbird is also to the right of the hawk. } \\
\\ 
\texttt{Next time, I will take a step back and re{-}evaluate the information. I will consider all the given statements and think about the relationships between the birds. I will also pay closer attention to the specific wording of the statements, such as "to the right of" and "to the left of". This will help me to avoid making the same mistake again. } \\
\\ 
\texttt{In this case, I should have noticed that the hummingbird is to the right of the hawk, which means that the cardinal cannot be the rightmost bird. The correct answer is actually (C) The hummingbird is the rightmost.} \\
            }
        }
    }
    & \\ \\

\end{supertabular}
}

\end{document}
